\scanpage{19}

associée à un $f\in F_M$, qui pour $f\in SF_M$ devient une bijection

\[
\alpha_f:\underline{\omega}_f\overset{\sim}{\longrightarrow}\{\pm e_M\}\tag*{(43)}
\]

d'où, grâce au choix de $e_M$, un élément

\[
\omega_f\in\underline{\omega}_f\qquad\alpha_f(\omega_f)=e_M\;,\tag*{(44)}
\]

\deleted{\unsure{qui}} i.e une \uline{orientation} \inserted{\uline{canonique}} de la face $f$.

La \deleted{\unsure{pro}}propriété de compatibilité (33) montre que les orientations des faces \inserted{adjacentes} sont compatibles deux à deux, donc définissent une orientation \inserted{$\mu$} de $SP_M$, appelée \uline{l'orientation canonique}.

On voit donc que le groupe $\mathrm{GL}_\Lambda(M)^\pm$ opère sur $SP_M$, et opère sur l'ens des orientations en \deleted{échangeant les \unsure{si|2}} \inserted{transformant l'orientation canonique} suivant $\det(g)$.

\[
g.\mu=\deleted{\unsure{\operatorname{sg}}}\det(g).\mu\qquad g\in\mathrm{GL}_\Lambda(M)^{\pm1}\tag*{(45).}
\]

Revenons \inserted{d'abord} à la question de la connexité

% archive sheet number: 18
